% Compile with: latexmk -pdf main
\documentclass[11pt]{article}

\usepackage[margin=1in]{geometry}
\usepackage{amsmath,amssymb,amsthm}
\usepackage[T1]{fontenc}
\usepackage[utf8]{inputenc}
\usepackage{lmodern}
\usepackage{newunicodechar}
\usepackage{fancyvrb}
\usepackage{xcolor}
\usepackage[hidelinks]{hyperref}
\usepackage{booktabs}

% Unicode characters appearing in Lean code
\newunicodechar{‖}{\ensuremath{\|}}
\newunicodechar{•}{\ensuremath{\bullet}}
\newunicodechar{₀}{\ensuremath{{}_0}}
\newunicodechar{₁}{\ensuremath{{}_1}}
\newunicodechar{₂}{\ensuremath{{}_2}}
\newunicodechar{ℝ}{\ensuremath{\mathbb{R}}}
\newunicodechar{ℂ}{\ensuremath{\mathbb{C}}}
\newunicodechar{ℕ}{\ensuremath{\mathbb{N}}}
\newunicodechar{𝔄}{\ensuremath{\mathfrak{A}}}
\newunicodechar{Λ}{\ensuremath{\Lambda}}
\newunicodechar{Φ}{\ensuremath{\Phi}}
\newunicodechar{κ}{\ensuremath{\kappa}}
\newunicodechar{γ}{\ensuremath{\gamma}}
\newunicodechar{ℋ}{\ensuremath{\mathcal{H}}}
\newunicodechar{→}{\ensuremath{\to}}
\newunicodechar{∀}{\ensuremath{\forall}}
\newunicodechar{≠}{\ensuremath{\neq}}
\newunicodechar{≤}{\ensuremath{\leq}}
\newunicodechar{∈}{\ensuremath{\in}}
\newunicodechar{∩}{\ensuremath{\cap}}
\newunicodechar{∑}{\ensuremath{\sum}}
\newunicodechar{∫}{\ensuremath{\int}}

\DefineVerbatimEnvironment{lean}{Verbatim}{fontsize=\small,frame=leftline,framerule=0.6pt,
  rulecolor=\color{gray!60},xleftmargin=1em}

\newtheorem{theorem}{Theorem}
\newtheorem{lemma}[theorem]{Lemma}
\theoremstyle{definition}
\newtheorem{definition}[theorem]{Definition}
\theoremstyle{remark}
\newtheorem{remark}[theorem]{Remark}

\newcommand{\fA}{\mathfrak{A}}
\newcommand{\norm}[1]{\|#1\|}
\newcommand{\Far}{\mathrm{Far}}

\title{A Formal Proof of the Lieb--Robinson Bound\\ and Exponential Clustering\\ on Finite Lattices in Lean~4}
\author{Nichlas Löfdahl \and Claude (Anthropic)\thanks{AI system (model Claude Opus 5.5), developed by Anthropic.
Claude wrote the Lean proofs and drafted this manuscript in collaboration with the first author.}}
\date{October 2026}

\begin{document}
\sloppy
\maketitle

\begin{abstract}
The Lieb--Robinson bound states that in a quantum lattice system with local interactions,
commutators of time-evolved local observables are exponentially small outside an effective light
cone, with a velocity independent of the size of the system. We present a machine-checked proof
of this bound for finite lattices in the Lean~4 proof assistant, on top of Mathlib. The setting
is a C$^*$-algebra with a subalgebra of observables for each region, subject only to the axiom that
observables in disjoint regions commute, so no tensor products are needed. For an interaction of
strength $J$, range $R$ and coordination $\kappa$ we prove
$\norm{[\tau_t(A),B]}\le 2e\,\norm{A}\norm{B}\,(n_X/\kappa)\,e^{v|t|-d(X,Y)/R}$ with $v=2eJ\kappa$,
where $n_X$ is the number of interaction terms meeting the support of $A$. The proof removes the
norm-preserving part of the dynamics with an explicit interaction-picture unitary, which avoids
time-dependent propagators, and iterates the resulting integral inequality by induction on
distance in the overlap graph of the interaction. We instantiate the abstract theorem for the
concrete lattice model of the Physlib library. As the main application we prove the exponential
clustering theorem: in a unique ground state with spectral gap $\gamma$, connected correlations
of observables at distance $d$ decay like $e^{-\mu d}$ with
$\mu=\min(1/2,\gamma/(8eJ\kappa))/R$. The spectral filter in its proof is evaluated with Fubini's
theorem and the Fourier transform of the Gaussian, without the error function or principal
values. The development uses no axioms beyond Lean's standard three.
\end{abstract}

\section{Introduction}

Lieb and Robinson showed in 1972 that quantum spin systems with finite-range interactions have a
finite group velocity \cite{liebrobinson1972}. Their bound became a central tool of quantum
many-body physics only three decades later, when Hastings, Koma, Nachtergaele, Ogata, Sims and
others used it to prove exponential clustering of correlations in gapped ground states
\cite{nachtergaelesims2006,hastingskoma2006}, to control the propagation of correlations
\cite{nachtergaeleogatasims2006}, and, among many other results, to prove the higher-dimensional
Lieb--Schultz--Mattis theorem and the stability of topological order. A survey is
\cite{nachtergaelesims2010}.

Despite this central role, we are not aware of a formalization of the Lieb--Robinson bound in any
proof assistant. Before starting we searched the web, the Physlib library
\cite{toobysmith2024heplean} and its open pull requests and issues. Physlib contains the
kinematics of quantum lattice systems (nets of local algebras, the quasi-local algebra,
interactions, finite-volume Hamiltonians and Gibbs states) but no time evolution and no locality
estimate. This note fills that gap for finite lattices.

\paragraph{Contributions.}
\begin{itemize}
  \item An abstract setting that uses only the commutation of observables in disjoint regions
  (\S\ref{sec:setting}).
  \item The basic integral inequality, proved with an explicit interaction-picture unitary
  instead of a time-dependent propagator (\S\ref{sec:onestep}).
  \item Its iteration by induction on a recursive notion of distance in the overlap graph of the
  interaction, without matrices or explicit chains (\S\ref{sec:iteration}).
  \item The bound in the familiar metric form $e^{v|t|-d/R}$ (\S\ref{sec:metric}), and its
  instantiation for Physlib's concrete lattice model (\S\ref{sec:lattice}).
  \item The exponential clustering theorem for gapped ground states (\S\ref{sec:clustering}).
\end{itemize}

The source is available in the \texttt{lieb-robinson} folder of
\url{https://github.com/Crotalus/ai-formalized}; it builds against Physlib commit \texttt{3bfc7af}
and Lean~\texttt{v4.35.0}.

\section{Setting}\label{sec:setting}

Let $\Lambda$ be a finite set of sites and $\fA$ a C$^*$-algebra. Physically $\fA$ is the algebra
of all observables of the system, for instance $\bigotimes_{x\in\Lambda}M_{n_x}(\mathbb C)$. The
usual construction through tensor products is cumbersome in current Mathlib, and the
Lieb--Robinson bound needs very little of it. A \emph{local net} assigns a subalgebra
$\fA_X\subseteq\fA$ to every region $X\subseteq\Lambda$, and we assume only that observables in
disjoint regions commute:

\begin{lean}
structure LocalNet (Λ 𝔄 : Type*) [CStarAlgebra 𝔄] where
  alg : Finset Λ → Subalgebra ℂ 𝔄
  commute_of_disjoint : ∀ {X Y : Finset Λ}, Disjoint X Y →
    ∀ a ∈ alg X, ∀ b ∈ alg Y, Commute a b

structure Interaction (N : LocalNet Λ 𝔄) where
  Φ : Finset Λ → 𝔄
  mem_alg : ∀ Z, Φ Z ∈ N.alg Z
  isSelfAdjoint : ∀ Z, IsSelfAdjoint (Φ Z)
\end{lean}

We do not even assume that $\fA_X\subseteq\fA_Y$ for $X\subseteq Y$: the proof never uses it. An
interaction assigns a self-adjoint term $\Phi(Z)\in\fA_Z$ to each region. Its Hamiltonian is
$H=\sum_Z\Phi(Z)$ and the Heisenberg dynamics is $\tau_t(A)=e^{itH}Ae^{-itH}$, defined with
Mathlib's exponential in Banach algebras. For self-adjoint $K$ the exponential $e^{itK}$ is
unitary, so $\tau_t$ is a $*$-automorphism that preserves norms. The commutator is
$[a,b]=ab-ba$.

The constants of the bound are:
\begin{itemize}
  \item the \emph{strength} $J$, with $\norm{\Phi(Z)}\le J$ for all $Z$;
  \item the \emph{coordination} $\kappa$: at most $\kappa$ regions $Z'$ with $\Phi(Z')\neq0$ meet
  any region $Z$ with $\Phi(Z)\ne0$ (Lean: \texttt{(I.neighbours Z).card ≤ κ});
  \item for the metric form, the \emph{range} $R$: every non-zero term has diameter at most $R$.
\end{itemize}

\section{The integral inequality}\label{sec:onestep}

\begin{lemma}\label{lem:onestep}
For $A\in\fA_X$, any $B\in\fA$ and $t\ge0$,
\[
  \norm{[\tau_t(A),B]}\le\norm{[A,B]}
  +2\norm{A}\sum_{Z\cap X\ne\emptyset}\int_0^t\norm{[\tau_s(\Phi(Z)),B]}\,ds .
\]
\end{lemma}

\begin{proof}
Split $H=H_1+H_2$, where $H_2$ collects the terms whose region meets $X$. Every term of $H_1$ lives
in a region disjoint from $X$, so $H_1$ commutes with $A$ and with $e^{isH_1}$. Put
$U(t)=e^{itH}e^{-itH_1}$ and $V(t)=e^{itH_1}e^{-itH}$, so that $UV=1$ and
$\tau_t(A)=U(t)\,A\,V(t)$. Then
\[
  [\tau_t(A),B]=U(t)\,[A,W(t)]\,V(t),\qquad W(t)=V(t)\,B\,U(t),
\]
and since $\norm{U},\norm{V}\le1$ we get
$\norm{[\tau_t(A),B]}\le\norm{[A,B]}+2\norm A\,\norm{W(t)-W(0)}$. A direct computation gives
\[
  W'(s)=e^{isH_1}\,\big[\tau_{-s}(B),\,iH_2\big]\,e^{-isH_1},
\]
whose norm is at most $\norm{[H_2,\tau_{-s}(B)]}=\norm{\tau_{-s}([\tau_s(H_2),B])}\le
\sum_{Z\cap X\neq\emptyset}\norm{[\tau_s(\Phi(Z)),B]}$. Integrating this bound on $W'$ gives
the claim.
\end{proof}

The textbook proof \cite{nachtergaelesims2006} obtains the same inequality from the
differential equation $g'(t)=i[\tau_t(H_2),g(t)]+\dots$ for $g(t)=[\tau_t(A),B]$, using that the
solution operator of $g'=i[K(t),g]$ is isometric. Formalizing this requires the time-dependent
propagator of that equation. The interaction-picture unitary $U$ achieves the same effect using
only exponentials of fixed operators. A Grönwall argument applied directly to $g$ would avoid
both but gives a weaker inequality, in which the supports of later terms grow with each step;
iterating it gives only a short-time bound.

In Lean the derivative of $W$ is computed with \texttt{HasDerivAt.mul} for the noncommutative
product of four factors and checked against the claimed form by \texttt{noncomm\_ring}. The
integrated bound is Mathlib's comparison theorem
\texttt{image\_norm\_le\_of\_norm\_deriv\_right\_le\_deriv\_boundary}, with the bounding function
$t\mapsto\int_0^t g$ differentiated by the fundamental theorem of calculus.

\begin{lean}
theorem norm_commutator_evolution_le {X : Finset Λ} {A : 𝔄} (hA : A ∈ N.alg X) (B : 𝔄)
    {t : ℝ} (ht : 0 ≤ t) :
    ‖commutator (I.evolution t A) B‖ ≤ ‖commutator A B‖ +
      2 * ‖A‖ * ∑ Z ∈ touching X, ∫ s in (0 : ℝ)..t,
        ‖commutator (I.evolution s (I.Φ Z)) B‖
\end{lean}

\section{Iteration}\label{sec:iteration}

Fix $B\in\fA_Y$ and write $f_Z(t)=\norm{[\tau_t(\Phi(Z)),B]}$. Lemma~\ref{lem:onestep} with
$A=\Phi(Z)$ gives $f_Z(t)\le f_Z(0)+2\norm{\Phi(Z)}\sum_{Z'\cap Z\ne\emptyset}\int_0^tf_{Z'}$, and
$f_Z(0)=0$ when $Z$ is disjoint from $Y$. Rather than expanding this into a sum over chains, we
measure distance to $Y$ in the overlap graph by a recursive predicate:

\begin{lean}
def Far (Y : Finset Λ) : ℕ → Finset Λ → Prop
  | 0, _ => True
  | k + 1, Z => I.Φ Z = 0 ∨ (Disjoint Z Y ∧ ∀ Z', (Z' ∩ Z).Nonempty → Far Y k Z')
\end{lean}

$\Far\,Y\,k\,Z$ says that every chain of overlapping regions with non-zero terms from $Z$ to a
region meeting $Y$ has at least $k$ steps. A single induction on $k$ then proves:

\begin{lemma}
If $\Far\,Y\,k\,Z$ and $t\ge0$, then
$f_Z(t)\le2\norm{\Phi(Z)}\norm B\,(2J\kappa t)^k/k!$.
\end{lemma}

The base case is $\norm{[a,b]}\le2\norm a\norm b$ together with $\norm{\tau_t(a)}=\norm{a}$. In the step,
the terms with $\Phi(Z')=0$ vanish, at most $\kappa$ terms remain, and
$\int_0^t(2J\kappa s)^k/k!\,ds=(2J\kappa)^kt^{k+1}/(k+1)!$. One more application of
Lemma~\ref{lem:onestep} gives the bound for a general $A$:

\begin{theorem}
Let $A\in\fA_X$ and $B\in\fA_Y$ with $X\cap Y=\emptyset$, and suppose every region meeting $X$
satisfies $\Far\,Y\,n$. Then for $t\ge0$
\[
  \norm{[\tau_t(A),B]}\le2\norm A\norm B\,n_X\,2J\,(2J\kappa)^n\,\frac{t^{n+1}}{(n+1)!}
  \le2\norm A\norm B\,\frac{n_X}{\kappa}\,e^{\,2eJ\kappa t-(n+1)} ,
\]
where $n_X$ is the number of non-zero terms meeting $X$ (the second form needs $\kappa>0$).
\end{theorem}

The exponential form uses $x^m/m!\le e^{ex-m}$, which follows from $(ex)^m/m!\le e^{ex}$. No
infinite series is needed: the polynomial bound is already exact at finite order.

\section{The metric form}\label{sec:metric}

Let $\Lambda$ carry a (pseudo)metric and let every non-zero term have diameter at most $R>0$. A
chain of $j+1$ overlapping terms spans a distance of at most $(j+1)R$, so a region all of whose
points are further than $kR$ from $Y$ is $\Far\,Y\,k$ (lemma \texttt{far\_of\_dist}). Choosing
$n=\lceil d/R\rceil-2$ (truncated at $0$) yields:

\begin{theorem}[Lieb--Robinson bound]
If every point of $X$ is at distance at least $d>0$ from every point of $Y$, then for
$A\in\fA_X$, $B\in\fA_Y$ and all real $t$,
\[
  \norm{[\tau_t(A),B]}\le 2e\,\norm A\norm B\,\frac{n_X}{\kappa}\,
  e^{\,v|t|-d/R},\qquad v=2eJ\kappa .
\]
\end{theorem}

\begin{lean}
theorem lieb_robinson_abs {X Y : Finset Λ} {A B : 𝔄} (hA : A ∈ N.alg X) (hB : B ∈ N.alg Y)
    {J R d : ℝ} {κ : ℕ} (hκ0 : 0 < κ) (hJ0 : 0 ≤ J) (hJ : ∀ Z, ‖I.Φ Z‖ ≤ J)
    (hκ : ∀ Z, (I.neighbours Z).card ≤ κ) (hR0 : 0 < R)
    (hR : ∀ Z, I.Φ Z ≠ 0 → ∀ x ∈ Z, ∀ y ∈ Z, dist x y ≤ R)
    (hd0 : 0 < d) (hd : ∀ x ∈ X, ∀ y ∈ Y, d ≤ dist x y) (t : ℝ) :
    ‖commutator (I.evolution t A) B‖ ≤
      2 * Real.exp 1 * ‖A‖ * ‖B‖ * ((I.neighbours X).card / κ) *
        Real.exp (2 * Real.exp 1 * J * κ * |t| - d / R)
\end{lean}

Negative times reduce to positive ones: the interaction $-\Phi$ satisfies the same hypotheses and
generates $\tau_{-t}$. None of the constants $J$, $R$, $\kappa$, $n_X$ depends on $|\Lambda|$.

\section{The concrete lattice model}\label{sec:lattice}

Physlib models a lattice system with a finite configuration space $\Omega_s$ at each site $s$ of a
possibly infinite set $S$. The operators of a finite region $X\subseteq S$ are the complex matrices
indexed by configurations of $X$, written \texttt{Alg Ω X}, and a larger region $Y\supseteq X$
contains them through $M\mapsto M\otimes1$ (\texttt{inclMat}). An interaction assigns a Hermitian
operator to each finite region, and the Hamiltonian of a finite box $\Lambda_0$ is the sum of the
lifts of the terms inside it. To apply the abstract theorem we fix a finite box $\Lambda_0$ and
take:
\begin{itemize}
  \item sites: the points of $\Lambda_0$, a finite type;
  \item algebra: \texttt{Alg Ω Λ₀}, a matrix algebra, with Mathlib's C$^*$-norm (the operator norm
  for the $\ell^2$ norm on configurations, \texttt{Matrix.Norms.L2Operator});
  \item local algebras: the images of the lifts $M\mapsto M\otimes1$ from each subregion.
\end{itemize}
The only property to verify is that lifts from disjoint regions commute. We prove this by
computing both products $(M\otimes1)(1\otimes N)$ and $(1\otimes N)(M\otimes1)$ entrywise: each
has a single non-zero term in the matrix product, at the configuration obtained by gluing.

\begin{lean}
lemma commute_inclMat_of_disjoint (hX : X ≤ Λ) (hY : Y ≤ Λ) (hd : Disjoint X Y)
    (M : Alg Ω X) (N : Alg Ω Y) : Commute (inclMat hX M) (inclMat hY N)
\end{lean}

As a check that the two settings fit together, the Hamiltonian of the abstract interaction
built from a Physlib interaction equals Physlib's Hamiltonian of the box
(\texttt{hamiltonian\_toBox}), after reindexing the sum over subsets of the box. The resulting
corollary states the Lieb--Robinson bound directly for Physlib's operators $A\in$ \texttt{Alg Ω X}
and $B\in$ \texttt{Alg Ω Y} with $X,Y\subseteq\Lambda_0$, the Physlib Hamiltonian, and strength
and range hypotheses on the Physlib interaction. Since the lift is contractive, the norms of $A$
and $B$ may be taken in their own regions. The coordination constant is stated for the box
interaction; formulating it on the infinite set of regions of $S$ is a matter of bookkeeping.


\section{Exponential clustering}\label{sec:clustering}

The most important consequence of the Lieb--Robinson bound is that a spectral gap forces
correlations to decay exponentially \cite{nachtergaelesims2006,hastingskoma2006}. We take
$\fA=\mathcal B(\mathcal H)$ for a finite-dimensional Hilbert space $\mathcal H$, which Mathlib
equips with its C$^*$-algebra structure.

\begin{definition}
A \emph{gapped ground state} of a self-adjoint $H$ with gap $\gamma>0$ is a unit vector $\Omega$
with $H\Omega=E_0\Omega$ such that $\langle\varphi,H\varphi\rangle\ge(E_0+\gamma)\norm{\varphi}^2$
for all $\varphi\perp\Omega$. The connected correlation is
$\langle AB\rangle-\langle A\rangle\langle B\rangle$, with
$\langle X\rangle=\langle\Omega,X\Omega\rangle$.
\end{definition}

The proof has three independent parts.

\paragraph{Spectral expansion.} Expanding in an orthonormal eigenbasis $\varphi_j$ of $H$, with
energies $E_0+\omega_j$, the function $F(t)=\langle[\tau_t(A),B]\rangle$ becomes
\[
  F(t)=F(0)+\sum_j c_j\,(e^{-i\omega_jt}-1)-\sum_j d_j\,(e^{i\omega_jt}-1),
\]
with $\sum_jc_j=\langle AB\rangle-\langle A\rangle\langle B\rangle$, $\sum_j|c_j|,\sum_j|d_j|\le
\norm A\norm B$ (Cauchy--Schwarz and Parseval), and only frequencies $\omega_j\ge\gamma$
contributing. The last point uses the gap: an eigenvector with energy below $E_0+\gamma$ lies in
the span of $\Omega$, so it does not contribute to the connected parts. To state the expansion
with a single frequency list we replace each $\omega_j$ by $\max(\omega_j,\gamma)$, which changes
only terms with zero coefficient. In Lean this needs the action of the exponential on
eigenvectors, $e^{K}v=e^{\mu}v$ for $Kv=\mu v$, which we prove from the exponential series.

\paragraph{A Gaussian filter.} The positive-frequency part $\sum_jc_j$ is extracted from $F$ by
integrating against $e^{-\alpha t^2}/(it)$. Since $F(0)=\langle[A,B]\rangle=0$ for disjoint
supports, it suffices to know
\[
  \int_{\mathbb R}e^{-\alpha t^2}\,\frac{e^{-i\omega t}-1}{it}\,dt
  =-\sqrt{\pi/\alpha}\int_0^\omega e^{-\sigma^2/4\alpha}\,d\sigma .
\]
We prove this by writing $(e^{-i\omega t}-1)/(it)=-\int_0^\omega e^{-i\sigma t}\,d\sigma$,
exchanging the integrals and using Mathlib's Fourier transform of the Gaussian. The right-hand
side is odd in $\omega$ and within $(2\sqrt{\pi\alpha}/\gamma)\,e^{-\gamma^2/4\alpha}$ of $-\pi$
for $\omega\ge\gamma$. This avoids the error function and principal-value integrals, neither of
which Mathlib provides. Combining with the expansion gives
$\int e^{-\alpha t^2}F(t)/(it)\,dt=-2\pi\sum_jc_j+\text{(error)}$.

\paragraph{Time-domain bound.} The same integral is bounded by
$\int e^{-\alpha t^2}|F(t)|/|t|\,dt$. For $|t|\le T$ the Lieb--Robinson bound in polynomial form
gives $|F(t)|\le K|t|^{n+1}$; for $|t|>T$ we use $|F|\le2\norm A\norm B$ and the Gaussian tail.
Choosing $T=(n+1)/(4eJ\kappa)$ and $\alpha=\gamma/(2T)$ balances the three errors and yields:

\begin{theorem}[Exponential clustering]
Let $\Omega$ be a gapped ground state of $H=\sum_Z\Phi(Z)$ with gap $\gamma$, let $A\in\fA_X$ and
$B\in\fA_Y$ with $X\cap Y=\emptyset$, and suppose every region meeting $X$ satisfies
$\Far\,Y\,n$. If $J,\kappa>0$, then
\[
  |\langle AB\rangle-\langle A\rangle\langle B\rangle|\le\frac2\pi\norm A\norm B
  \Big(\frac{n_X}{\kappa}e^{-(n+1)/2}+\Big(\frac{4eJ\kappa}\gamma
  +\sqrt{\frac{2\pi eJ\kappa}\gamma}\Big)e^{-\gamma(n+1)/(8eJ\kappa)}\Big).
\]
On a metric lattice with range $R$ and $d(X,Y)\ge d$, the right-hand side becomes
exponentially small in $d$ at rate $\min(1/2,\gamma/(8eJ\kappa))/R$.
\end{theorem}

\begin{lean}
theorem exponential_clustering {γ : ℝ} (hγ : 0 < γ)
    (G : GappedGroundState I.hamiltonian γ) {X Y : Finset Λ} {A B : ℋ →L[ℂ] ℋ}
    (hA : A ∈ N.alg X) (hB : B ∈ N.alg Y) (hXY : Disjoint X Y) {J : ℝ} {κ : ℕ}
    (hκ0 : 0 < κ) (hJ0 : 0 < J) (hJ : ∀ Z, ‖I.Φ Z‖ ≤ J)
    (hκ : ∀ Z, (I.neighbours Z).card ≤ κ) {n : ℕ}
    (hn : ∀ Z, (Z ∩ X).Nonempty → I.Far Y n Z) :
    ‖G.corr A B‖ ≤ 2 / Real.pi * ‖A‖ * ‖B‖ *
      ((I.neighbours X).card / κ * Real.exp (-(n + 1) / 2) +
        (4 * Real.exp 1 * J * κ / γ + Real.sqrt (2 * Real.pi * Real.exp 1 * J * κ / γ)) *
          Real.exp (-γ * (n + 1) / (8 * Real.exp 1 * J * κ)))
\end{lean}

The rate $\gamma/(8eJ\kappa)\sim\gamma/v$ is of the form $\gamma/v$ up to constants, as in the
rates obtained by Nachtergaele and Sims and by Hastings and Koma
\cite{nachtergaelesims2006,hastingskoma2006}; we did not optimize constants. We assume a unique ground state. The degenerate case,
where correlations are taken in a ground state projection, is not treated.

\section{Discussion}\label{sec:discussion}

\paragraph{Size and checking.} The development consists of 18 definitions and 85 lemmas in
about 2200 lines, including documentation, in nine files: the setting (\S\ref{sec:setting}), the
integral inequality (\S\ref{sec:onestep}), the iteration (\S\ref{sec:iteration}), the metric form
(\S\ref{sec:metric}), the lattice model (\S\ref{sec:lattice}), and four files for exponential
clustering (\S\ref{sec:clustering}). \texttt{\#print axioms} on the
main results reports only \texttt{propext}, \texttt{Classical.choice} and \texttt{Quot.sound}. We
have not run Physlib's contribution linters on these files.

\paragraph{Design choices.} Three decisions kept the Lieb--Robinson part short. First, the axiom that
disjoint regions commute replaces tensor products and isotony of the net; the bound uses nothing
else about locality. Second, the interaction-picture unitary replaces the time-dependent
propagator of the textbook proof. Third, the recursive predicate $\Far$ replaces the sum over
chains: each induction step needs only one application of the integral inequality and a count
of neighbours. The proof was developed by first fixing all statements and checking them with
\texttt{sorry}, after which the files were proved in parallel by separate instances of Claude
against the frozen interface, and the lattice model was added in the same way. Every claim was then re-verified by a full build.

\paragraph{Constants.} Our velocity $v=2eJ\kappa$ has the same structure as in the literature
\cite{nachtergaelesims2006,nachtergaelesims2010}, where $v$ is proportional to the interaction
strength and a local summability constant. The factor $e$ comes from the elementary estimate
$x^m/m!\le e^{ex-m}$. The prefactor depends on the number $n_X$ of terms meeting $X$, as is
standard. We did not optimize constants.

\paragraph{Limitations and future work.} The result is for finite lattices. The bounds are uniform
in the volume, which is what the infinite-volume theory needs, but we do not construct the
infinite-volume dynamics. Interactions are of finite range: exponentially decaying interactions
would require the summability (``$F$-function'') formalism of \cite{nachtergaeleogatasims2006}.
Exponential clustering is proved for a unique gapped ground state in finite dimension; the
degenerate case and the infinite-volume formulation, where the gap is a property of a
thermodynamic limit, are not covered. Further applications such as area laws or the
Lieb--Schultz--Mattis theorem would require substantially more infrastructure.

\paragraph{Author contributions.} N.L.\ chose the problem, directed the work and reviewed the
results. Claude, an AI system developed by Anthropic, surveyed the literature and libraries,
designed the statements and the proof architecture, wrote the Lean proofs (in part through
parallel sub-agents) and drafted the manuscript. All statements were checked by the Lean kernel.
As an AI system cannot take responsibility for a publication, N.L.\ is accountable for its
content.

\bibliographystyle{plain}
\bibliography{refs}

\end{document}
